% 图 2-1 顶层架构 v2 — 修：删 ①②③ 流程编号（架构图非时序图）/ 重排 3 条 agent→bus 箭头错开 endpoint
\documentclass[border=4pt]{standalone}
\usepackage{xeCJK}
\setCJKmainfont{Songti SC}
\setCJKsansfont{Heiti SC}
\usepackage{tikz}
\usepackage{amsmath}
\usetikzlibrary{arrows.meta, positioning, calc, shapes.geometric, fit, backgrounds}
% 共享 TikZ 风格 — Aegaeon (SOSP'25) inspired，对标顶刊系统论文
% 使用：% 共享 TikZ 风格 — Aegaeon (SOSP'25) inspired，对标顶刊系统论文
% 使用：% 共享 TikZ 风格 — Aegaeon (SOSP'25) inspired，对标顶刊系统论文
% 使用：\input{_style.tex} 在 standalone preamble 内
% 调色：Okabe-Ito（与 matplotlib 图一致）

\definecolor{fuxiBlue}{HTML}{0072B2}    % 自家系统主色
\definecolor{fuxiOrange}{HTML}{E69F00}  % 第一对照
\definecolor{fuxiGreen}{HTML}{009E73}
\definecolor{fuxiPurple}{HTML}{CC79A7}
\definecolor{fuxiSky}{HTML}{56B4E9}
\definecolor{fuxiGrayBox}{HTML}{F2F2F2} % 浅灰底
\definecolor{fuxiGrayLine}{HTML}{808080}
\definecolor{fuxiGrayDark}{HTML}{555555}

\tikzset{
  % ── 通用节点 ────────────────────────────────────────
  modBox/.style = {
    draw=fuxiGrayLine, thick, fill=white,
    rounded corners=2pt, inner sep=4pt,
    minimum height=8mm, minimum width=20mm,
    align=center, font=\small
  },
  highlightBox/.style = {
    modBox, draw=fuxiBlue, fill=fuxiBlue!8,
    text=fuxiBlue, font=\small\bfseries
  },
  ghostBox/.style = {
    modBox, draw=fuxiGrayLine!50, fill=fuxiGrayBox,
    text=fuxiGrayDark, font=\small\itshape
  },
  layerBg/.style = {
    draw=fuxiGrayLine!40, dashed, fill=fuxiGrayBox!40,
    rounded corners=4pt, inner sep=8pt
  },
  externalBoundary/.style = {
    draw=fuxiGrayLine, dashed, dash pattern=on 4pt off 3pt,
    rounded corners=6pt, inner sep=10pt
  },
  % ── 箭头 ─────────────────────────────────────────────
  flow/.style = {
    -{Stealth[length=4pt, width=4pt]},
    line width=0.7pt, color=fuxiGrayDark
  },
  flowEvent/.style = {
    flow, color=fuxiBlue, line width=0.9pt
  },
  flowDup/.style = {
    flow, color=fuxiBlue!70, line width=0.7pt, dotted
  },
  bidir/.style = {
    {Stealth[length=4pt]}-{Stealth[length=4pt]},
    line width=0.7pt, color=fuxiGrayDark
  },
  % ── 编号气泡 ① ② ③ ────────────────────────────────
  numBubble/.style = {
    circle, draw=fuxiBlue, fill=white, text=fuxiBlue,
    font=\scriptsize\bfseries, inner sep=0pt,
    minimum size=4mm, line width=0.6pt
  },
}
 在 standalone preamble 内
% 调色：Okabe-Ito（与 matplotlib 图一致）

\definecolor{fuxiBlue}{HTML}{0072B2}    % 自家系统主色
\definecolor{fuxiOrange}{HTML}{E69F00}  % 第一对照
\definecolor{fuxiGreen}{HTML}{009E73}
\definecolor{fuxiPurple}{HTML}{CC79A7}
\definecolor{fuxiSky}{HTML}{56B4E9}
\definecolor{fuxiGrayBox}{HTML}{F2F2F2} % 浅灰底
\definecolor{fuxiGrayLine}{HTML}{808080}
\definecolor{fuxiGrayDark}{HTML}{555555}

\tikzset{
  % ── 通用节点 ────────────────────────────────────────
  modBox/.style = {
    draw=fuxiGrayLine, thick, fill=white,
    rounded corners=2pt, inner sep=4pt,
    minimum height=8mm, minimum width=20mm,
    align=center, font=\small
  },
  highlightBox/.style = {
    modBox, draw=fuxiBlue, fill=fuxiBlue!8,
    text=fuxiBlue, font=\small\bfseries
  },
  ghostBox/.style = {
    modBox, draw=fuxiGrayLine!50, fill=fuxiGrayBox,
    text=fuxiGrayDark, font=\small\itshape
  },
  layerBg/.style = {
    draw=fuxiGrayLine!40, dashed, fill=fuxiGrayBox!40,
    rounded corners=4pt, inner sep=8pt
  },
  externalBoundary/.style = {
    draw=fuxiGrayLine, dashed, dash pattern=on 4pt off 3pt,
    rounded corners=6pt, inner sep=10pt
  },
  % ── 箭头 ─────────────────────────────────────────────
  flow/.style = {
    -{Stealth[length=4pt, width=4pt]},
    line width=0.7pt, color=fuxiGrayDark
  },
  flowEvent/.style = {
    flow, color=fuxiBlue, line width=0.9pt
  },
  flowDup/.style = {
    flow, color=fuxiBlue!70, line width=0.7pt, dotted
  },
  bidir/.style = {
    {Stealth[length=4pt]}-{Stealth[length=4pt]},
    line width=0.7pt, color=fuxiGrayDark
  },
  % ── 编号气泡 ① ② ③ ────────────────────────────────
  numBubble/.style = {
    circle, draw=fuxiBlue, fill=white, text=fuxiBlue,
    font=\scriptsize\bfseries, inner sep=0pt,
    minimum size=4mm, line width=0.6pt
  },
}
 在 standalone preamble 内
% 调色：Okabe-Ito（与 matplotlib 图一致）

\definecolor{fuxiBlue}{HTML}{0072B2}    % 自家系统主色
\definecolor{fuxiOrange}{HTML}{E69F00}  % 第一对照
\definecolor{fuxiGreen}{HTML}{009E73}
\definecolor{fuxiPurple}{HTML}{CC79A7}
\definecolor{fuxiSky}{HTML}{56B4E9}
\definecolor{fuxiGrayBox}{HTML}{F2F2F2} % 浅灰底
\definecolor{fuxiGrayLine}{HTML}{808080}
\definecolor{fuxiGrayDark}{HTML}{555555}

\tikzset{
  % ── 通用节点 ────────────────────────────────────────
  modBox/.style = {
    draw=fuxiGrayLine, thick, fill=white,
    rounded corners=2pt, inner sep=4pt,
    minimum height=8mm, minimum width=20mm,
    align=center, font=\small
  },
  highlightBox/.style = {
    modBox, draw=fuxiBlue, fill=fuxiBlue!8,
    text=fuxiBlue, font=\small\bfseries
  },
  ghostBox/.style = {
    modBox, draw=fuxiGrayLine!50, fill=fuxiGrayBox,
    text=fuxiGrayDark, font=\small\itshape
  },
  layerBg/.style = {
    draw=fuxiGrayLine!40, dashed, fill=fuxiGrayBox!40,
    rounded corners=4pt, inner sep=8pt
  },
  externalBoundary/.style = {
    draw=fuxiGrayLine, dashed, dash pattern=on 4pt off 3pt,
    rounded corners=6pt, inner sep=10pt
  },
  % ── 箭头 ─────────────────────────────────────────────
  flow/.style = {
    -{Stealth[length=4pt, width=4pt]},
    line width=0.7pt, color=fuxiGrayDark
  },
  flowEvent/.style = {
    flow, color=fuxiBlue, line width=0.9pt
  },
  flowDup/.style = {
    flow, color=fuxiBlue!70, line width=0.7pt, dotted
  },
  bidir/.style = {
    {Stealth[length=4pt]}-{Stealth[length=4pt]},
    line width=0.7pt, color=fuxiGrayDark
  },
  % ── 编号气泡 ① ② ③ ────────────────────────────────
  numBubble/.style = {
    circle, draw=fuxiBlue, fill=white, text=fuxiBlue,
    font=\scriptsize\bfseries, inner sep=0pt,
    minimum size=4mm, line width=0.6pt
  },
}


\begin{document}
\begin{tikzpicture}[node distance=8mm and 14mm, font=\sffamily]

% ── 用户面（最左）───────────────────────────────────────
\node[modBox, minimum width=24mm] (im) {IM/PWA\\\scriptsize 浏览器/手机};
\node[modBox, below=8mm of im, minimum width=24mm] (tui) {Firehose TUI\\\scriptsize 终端实时观察};

% ── 编排层核心（中间，三层）──────────────────────────
\node[highlightBox, right=24mm of im, minimum width=30mm] (xuannv) {玄女 (xuannv)\\\scriptsize 顶层 agent};
\node[modBox, below=8mm of xuannv, minimum width=30mm] (shelf) {Shelf 橱柜\\\scriptsize 门客注册表};
\node[highlightBox, below=8mm of shelf, minimum width=30mm] (bus) {EventBus\\\scriptsize broadcast + WAL};

% ── 门客（最右）─────────────────────────────────────
\node[modBox, right=24mm of xuannv, minimum width=26mm] (cc1) {cc 门客 \#1\\\scriptsize claude-code};
\node[modBox, below=4mm of cc1, minimum width=26mm] (cc2) {cc 门客 \#2\\\scriptsize claude-code};
\node[modBox, below=4mm of cc2, minimum width=26mm] (codex) {codex 门客\\\scriptsize OpenAI codex};

% ── 持久化层（最下，居中）─────────────────────────
\node[ghostBox, below=10mm of bus, minimum width=30mm] (sqlite) {SQLite WAL\\\scriptsize 单一真相源};

% ── 跨进程边界（虚线）─────────────────────────────
\begin{scope}[on background layer]
  \node[externalBoundary, fit=(cc1) (cc2) (codex), inner sep=5pt,
    label={[font=\scriptsize\itshape, text=fuxiGrayDark, yshift=1pt]above:{子进程（CLI 头无 GUI）}}] (subprocs) {};
  \node[externalBoundary, fit=(im) (tui), inner sep=5pt,
    label={[font=\scriptsize\itshape, text=fuxiGrayDark, yshift=1pt]above:{用户终端/浏览器}}] (clients) {};
\end{scope}

% ── 控制流（灰，bidir）────────────────────────────
% 用户 → 玄女
\draw[flow] (im.east) -- node[above=0.5pt, font=\scriptsize] {REST/WS} (xuannv.west);

% 玄女 → 门客（A2A，bidir 但只一条主线，避免分叉乱）
\draw[bidir] (xuannv.east) -- node[above=0.5pt, font=\scriptsize] {A2A} (cc1.west);
\draw[bidir] (xuannv.east) -| ($(cc2.west)+(-5mm,0)$) -- (cc2.west);
\draw[bidir] (xuannv.east) -| ($(codex.west)+(-5mm,0)$) -- (codex.west);

% 玄女 ↔ Shelf, Shelf ↔ EventBus
\draw[bidir] (xuannv.south) -- (shelf.north);
\draw[bidir] (shelf.south) -- (bus.north);

% ── 事件流（蓝）：3 条 agent → bus 错开 endpoint，避免堆叠 ──
% cc1 进 bus 上沿
\draw[flowEvent] (cc1.west) -| ($(bus.east)+(8mm,0)$) |- ($(bus.east)+(0,4mm)$);
% cc2 进 bus 中
\draw[flowEvent] (cc2.west) -| ($(bus.east)+(6mm,0)$) |- (bus.east);
% codex 进 bus 下沿
\draw[flowEvent] (codex.west) -| ($(bus.east)+(4mm,0)$) |- ($(bus.east)+(0,-4mm)$);

% Firehose TUI → bus（订阅）
\draw[flow] (tui.east) -| ($(bus.west)+(-3mm,0)$) -- (bus.west);

% bus → SQLite（持久化）
\draw[flowEvent] (bus.south) -- (sqlite.north);

% bus → IM/PWA（loopback 订阅，浅蓝点线）
\draw[flowDup] (bus.west) -| ($(im.south)+(-4mm,0)$) -- ($(im.south)+(-4mm,3mm)$);

% ── 图例（右下角，紧凑）──────────────────────────
\node[anchor=north west, font=\scriptsize, align=left, draw=fuxiGrayLine!40, fill=white,
      inner sep=4pt, line width=0.3pt] (legend) at ($(sqlite.south east)+(8mm,0)$) {
\textcolor{fuxiGrayDark}{$\rightarrow$ 控制流 \quad}
\textcolor{fuxiBlue}{$\rightarrow$ 事件流（publish）\quad}
\textcolor{fuxiBlue!70}{$\cdots\!\rightarrow$ 订阅广播}
};

\end{tikzpicture}
\end{document}
